\SourcePage{364}{367}
\section{Electronic terms of a diatomic molecule}

The great disparity between atomic nuclear masses and the electron mass is fundamental to molecular theory. Owing to this mass difference, nuclei in a molecule move slowly in comparison with the electrons. We may therefore study the electronic motion with the nuclei held fixed at prescribed separations. The energy levels $U_n$ found for this system are called the molecule's \emph{electronic terms}. For atoms, the energy levels were definite numbers; here, by contrast, the electronic terms are functions of parameters---the separations of the nuclei in the molecule. The quantity $U_n$ also includes the electrostatic interaction energy of the nuclei. Thus $U_n$ is, in substance, the total energy of the molecule for a specified configuration of fixed nuclei.

We begin the study of molecules with the simplest type, the \emph{diatomic} molecule, for which the fullest theoretical investigation is possible. Its electronic terms depend on just one parameter: the internuclear distance $r$.

One of the principal ways of classifying atomic terms used the values of the total orbital angular momentum $L$. For molecules, however, the total orbital angular momentum of the electrons is not conserved in general, since the electric field of several nuclei has no central symmetry.

In a diatomic molecule the field does, nevertheless, possess axial symmetry about the line through both nuclei. The projection of the orbital angular momentum on this axis is consequently conserved, and molecular electronic terms can be classified by its values. The magnitude of this projection on the molecular axis is conventionally denoted by $\Lambda$; its possible values are $0,1,2,\ldots$ Terms with different $\Lambda$ are designated by the capital Greek letters corresponding to the Latin symbols for atomic

\SourcePage{365}{368}
terms having different $L$. Thus the cases $\Lambda=0,1,2$ are called, respectively, $\Sigma$-, $\Pi$-, and $\Delta$-terms; larger values of $\Lambda$ ordinarily need not be considered.

Every electronic state of a molecule is further characterized by the total spin $S$ of all its electrons. When $S$ is nonzero, the different orientations of the total spin give a degeneracy of multiplicity $2S+1$\footnote{We disregard here the fine structure arising from relativistic interactions (see below, \S~83, 84).}. As for atoms, the number $2S+1$ is called the term's \emph{multiplicity} and is placed to the upper left of the term symbol. For example, ${}^3\Pi$ denotes a term with $\Lambda=1$, $S=1$.

Besides rotations through an arbitrary angle about the axis, the symmetry of the molecule permits reflection in any plane containing that axis. Such a reflection leaves the molecular energy unchanged, although the resulting state need not be wholly identical to the initial one. Under reflection in a plane through the molecular axis, the sign of the angular momentum about that axis changes (angular momentum is an axial vector). It follows that every electronic term with nonzero $\Lambda$ is doubly degenerate: each energy belongs to two states that differ in the direction of the orbital-angular-momentum projection on the molecular axis. When $\Lambda=0$, on the other hand, reflection does not alter the molecular state, so $\Sigma$-terms are nondegenerate. Reflection can only multiply the wave function of a $\Sigma$-term by a constant. Since two reflections in the same plane amount to the identity transformation, this constant is $\pm1$. We must therefore distinguish the $\Sigma$-terms whose wave function remains entirely unchanged under reflection from those whose wave function reverses sign. The former are denoted by $\Sigma^+$ and the latter by $\Sigma^-$.

For a molecule composed of two identical atoms, a further symmetry arises and supplies an additional characteristic of its electronic terms. A diatomic molecule with identical nuclei also has a center of symmetry at the midpoint of the line joining them (we choose this point as the coordinate origin)\footnote{It also has a symmetry plane perpendicular to the molecular axis and bisecting it. There is no need to treat this symmetry element separately, since the existence of such a plane follows automatically from the presence of a center and an axis of symmetry.}. The Hamiltonian is therefore invariant under simultaneous reversal of the coordinates of all the molecular electrons, with the nuclear co-

\SourcePage{366}{369}
ordinates left unchanged. The operator of this transformation\footnote{This must not be confused with inversion of the coordinates of every particle in the system (compare \S~86)!} also commutes with the orbital-angular-momentum operator. Hence terms of definite $\Lambda$ can be classified additionally by parity: the wave function of an \emph{even} ($g$) state is unchanged when the electron coordinates reverse sign, whereas that of an \emph{odd} ($u$) state changes sign. The parity indices $u,g$ are conventionally written as subscripts on the term symbol: $\Pi_u$, $\Pi_g$, and so forth.

Empirically, the normal electronic state of the overwhelming majority of chemically stable diatomic molecules is known to be totally symmetric: the electronic wave function is invariant under every symmetry transformation of the molecule. In the great majority of cases, the total spin $S$ also vanishes in the normal state. In other words, the molecular ground term is ${}^1\Sigma^+$, or ${}^1\Sigma_g^+$ when the molecule consists of identical atoms. Familiar exceptions to these rules are the molecules $\mathrm O_2$ (normal term ${}^3\Sigma_g^-$) and $\mathrm{NO}$ (normal term ${}^2\Pi$).

\ProblemHeading Separate the variables in the Schrödinger equation for the electronic terms of the $\mathrm H_2^+$ ion, using elliptic coordinates.

\SolutionHeading The Schrödinger equation for an electron in the field of two fixed protons is
\[
 \frac{\Delta\psi}{2}+\left(E+\frac1{r_1}+\frac1{r_2}\right)\psi=0
\]
(atomic units are used). Elliptic coordinates are defined by
\[
 \xi=\frac{r_1+r_2}{R},\qquad \eta=\frac{r_2-r_1}{R},\qquad
 1\leqslant\xi<\infty,\quad -1\leqslant\eta\leqslant1,
\]
and the third coordinate is the angle $\varphi$ of rotation about the axis through the two nuclei, whose separation is $R$ (see I, \S~48). In these coordinates the Laplacian is
\[
 \Delta=\frac4{R^2(\xi^2-\eta^2)}\left[
 \frac{\partial}{\partial\xi}\left[(\xi^2-1)\frac{\partial}{\partial\xi}\right]
 +\frac{\partial}{\partial\eta}\left[(1-\eta^2)\frac{\partial}{\partial\eta}\right]
 \right]+\frac1{R^2(\xi^2-1)(1-\eta^2)}\frac{\partial^2}{\partial\varphi^2}.
\]
Setting
\[
 \psi=X(\xi)Y(\eta)e^{i\Lambda\varphi},
\]
we obtain the following equations for $X$ and $Y$:
\[
 \frac d{d\xi}\left[(\xi^2-1)\frac{dX}{d\xi}\right]
 +\left(\frac{ER^2\xi^2}{2}+2R\xi+A-\frac{\Lambda^2}{\xi^2-1}\right)X=0,
\]
\[
 \frac d{d\eta}\left[(1-\eta^2)\frac{dY}{d\eta}\right]
 +\left(-\frac{ER^2\eta^2}{2}-A-\frac{\Lambda^2}{1-\eta^2}\right)Y=0,
\]

\SourcePage{367}{370}
where $A$ is the separation parameter. Every electronic term $E(R)$ is characterized by three quantum numbers: $\Lambda$ and two ``elliptic quantum numbers'' $n_\xi$, $n_\eta$, which specify the numbers of zeros of the functions $X(\xi)$ and $Y(\eta)$.
